\documentclass[12pt]{article}

\usepackage[a4paper, margin=1in]{geometry}
\usepackage{amsmath, amssymb}
\usepackage{setspace}
\usepackage{hyperref}

\setstretch{1.2}

\title{\textbf{Developmental Brain GNN: \\ Emergent Functional Specialization and Phase Transitions in Growing Neural Graphs}}
\author{Iyer Ramkumar Ramasubramanian \\ \textit{Independent Researcher}}
\date{\today}

\begin{document}

\maketitle

\begin{abstract}
Understanding the emergence of modular brain structures remains a fundamental challenge in both neurobiology and artificial intelligence. This paper introduces the \textbf{Developmental Brain GNN}, a memory-safe graph growth model capable of simulating functional specialization from a single undifferentiated node. By tracking region-specific metrics over time, we demonstrate a phase transition from early global regulation (Hypothalamus dominance) to specialized sequence processing (Temporal dominance), ultimately highlighting the necessity of competing pressures for maintaining functional diversity.
\end{abstract}

\section{Introduction}
As an independent researcher investigating the convergence of artificial intelligence and neurobiology, I propose that brain-like structures can grow dynamically rather than being statically defined. The Developmental Brain GNN extends the "Homo Deus" framework by moving beyond a fixed architecture into a self-organizing system that matures through functional specialization.

\section{System Features and Architecture}

The Developmental Brain GNN implements a global workspace architecture with the following core components:

\begin{itemize}
\item \textbf{Memory-Safe Graph Growth}: The network expands from a single proto-cortical node into a complex multi-node graph.
\item \textbf{Functional Specialization}: Regions such as the Occipital (vision), Temporal (sequence), and Prefrontal (control) emerge dynamically based on input signals.
\item \textbf{Attention Hub}: A simple global workspace facilitates communication across specialized regions.
\end{itemize}

\section{Developmental Phases}

Experimental logs reveal a distinct three-stage developmental trajectory:

\subsection{Initial Undifferentiated State (Step 0)}
The system begins with a single node in a symmetric, undifferentiated proto-cortical state (Parietal region).

\subsection{Global Regulatory Phase (Step 50)}
As the network grows, it enters a phase dominated by the Hypothalamus, representing internal homeostatic regulation that precedes higher cognitive functions.

\subsection{Phase Transition and Specialization (Steps 100--150)}
The system undergoes a sharp transition where sequence-processing signals become dominant, leading to the rapid expansion of the Temporal region and eventual functional collapse into a stabilized monoculture.

\section{Mathematical and Dynamical Interpretation}

The network acts as a self-organizing system with strong attractors. Without diversity regularization or spatial constraints, the system naturally follows a "winner-take-all" specialization path:

\[
\text{Region}_{dominant} = \arg\max_{r} (\text{Signal Strength}_r)
\]

This highlights a fundamental principle: true brain structure requires balanced competing pressures to maintain modular diversity and prevent functional collapse.

\section{Results and Inference}

Observations from the simulation logs show:
\begin{itemize}
\item \textbf{Dynamic Growth}: A complex brain-like structure successfully grows from a single node.
\item \textbf{Emergent Attractors}: The Temporal region acts as a strong attractor in the current configuration, dominating the final network state.
\item \textbf{Constraint Requirement}: Future improvements must integrate spatial embedding and multi-objective optimization to achieve realistic modularity.
\end{itemize}

\section{Conclusion}
The Developmental Brain GNN demonstrates that functional specialization can emerge dynamically in growing neural graphs. By modeling the transition from homeostasis to cognitive specialization, we provide a framework for future research into adaptive brain architectures and distributed intelligence.

\section{Repository for Experimentation}
\noindent
\url{https://github.com/spacetracker-collab/BrainDevelopment}

\end{document}
